% auto-generated by build_deck.py — do not edit by hand
% All numerical values sourced from deck_values.json
\documentclass[aspectratio=169,10pt]{beamer}

\usetheme{default}
\usepackage[utf8]{inputenc}
\usepackage[T1]{fontenc}
\usepackage{booktabs}
\usepackage{graphicx}
\usepackage{amsmath}
\usepackage{xcolor}
\usepackage{siunitx}

\definecolor{shipblue}{RGB}{20,74,130}
\definecolor{warnorange}{RGB}{210,105,30}
\setbeamercolor{structure}{fg=shipblue}
\setbeamercolor{alerted text}{fg=warnorange}
\setbeamertemplate{navigation symbols}{}
\setbeamertemplate{footline}[frame number]

% ── macros sourced from deck_values.json ──────────────────────────────────────
\newcommand{\BarLen}{1400}
\newcommand{\NEv}{2,000}
\newcommand{\NPos}{31}
\newcommand{\LamBulk}{160}
\newcommand{\LamEffL}{\ensuremath{32.2\pm0.0}\,cm}
\newcommand{\LamEffR}{\ensuremath{31.4\pm0.0}\,cm}
\newcommand{\LamEffC}{\ensuremath{31.8\pm0.0}\,cm}
\newcommand{\MylarR}{0.90}
\newcommand{\MylarSL}{1.0}
\newcommand{\MylarSA}{0.1}
\newcommand{\PDEpeak}{63\%}
\newcommand{\SiPMModel}{AFBR-S4N66P014M}
\newcommand{\SiPMModelPDE}{AFBR-S4N66P024M}
\newcommand{\ScintYield}{10,000\,$\gamma$/MeV}
\newcommand{\TauRise}{0.7\,ns}
\newcommand{\TauDecay}{1.8\,ns}
\newcommand{\SigmaAtMsixnine}{$406\pm27$\,ps}
\newcommand{\SigmaAtMfour}{$211\pm11$\,ps}
\newcommand{\SigmaAtZero}{$142\pm4$\,ps}
\newcommand{\SigmaAtFour}{$209\pm9$\,ps}
\newcommand{\SigmaAtSixnine}{$372\pm19$\,ps}
\newcommand{\VeffFPT}{264.1\,mm/ns}
\newcommand{\VeffSUMfour}{213.7\,mm/ns}
\newcommand{\VeffTfifty}{275.7\,mm/ns}
\newcommand{\BudgetEsc}{4.84\%}
\newcommand{\BudgetSiPM}{5.77\%}
\newcommand{\BudgetDet}{3.63\%}
\newcommand{\NpeLeftZero}{77.2}
\newcommand{\NpeRightZero}{77.0}
\newcommand{\NpeLeftNear}{2926}
\newcommand{\NpeRightNear}{2914}

\title[End-only + Mylar EJ-204]{%
  End-only + Mylar-wrapped EJ-204 bar:\\
  Optical simulation study}
\subtitle{Geant4 11.4.0 / SSLG4-OPSim}
\author{R.\ R\'ios (ULS)}
\institute{SHiP Collaboration Meeting}
\date{2026-06-14}

\begin{document}

% ─── FRAME 1: Title ───────────────────────────────────────────────────────────
\begin{frame}
  \titlepage
\end{frame}

% ─── FRAME 2: Motivation & scope ─────────────────────────────────────────────
\begin{frame}{Motivation \& scope}
  \begin{columns}[T]
    \column{0.55\textwidth}
    \begin{itemize}
      \item \textbf{Deliberate simplification:} readout \emph{only} at the
            two ends ($\pm X$ faces); all other faces wrapped in Mylar.
      \item \textbf{Goals:}
        \begin{itemize}
          \item Isolate end-readout contribution.
          \item Verify exponential attenuation with Mylar recirculation.
          \item Benchmark the SUM4 timing estimator on a clean geometry.
        \end{itemize}
      \item \NPos\ beam positions ($x=-690\ldots+690$\,mm),
            \NEv\ $\mu^-$ events per position.
    \end{itemize}
    \column{0.43\textwidth}
    \begin{block}{This study}
      \textbf{End-only + Mylar} = extreme limit.
      Photons either exit to an end SiPM or are absorbed by Mylar
      after multiple reflections.
    \end{block}
    \begin{alertblock}{\small No TB counterpart}
      \small This configuration has \alert{no direct test-beam counterpart}
      (April-2026 TB used top readout + Tyvek).
    \end{alertblock}
    \begin{alertblock}{\small Electronic jitter not included}
      \small SPTR $\oplus$ FastIC jitter are post-processing steps,
      not folded into $\sigma_t$ shown here.
    \end{alertblock}
  \end{columns}
\end{frame}

% ─── FRAME 3: Geometry & simulation setup ────────────────────────────────────
\begin{frame}{Geometry \& simulation setup}
  \begin{columns}[T]
    \column{0.55\textwidth}
    \begin{itemize}
      \item \textbf{Bar:} EJ-204 (OPSC-101),
            $\BarLen \times 60 \times 10$\,mm$^3$.
      \item \textbf{Longitudinal axis:} $X$ ($-700$ to $+700$\,mm).
      \item \textbf{End-left SiPMs} (IDs 0--7): $-X$ face.
            \textbf{End-right SiPMs} (IDs 8--15): $+X$ face.
      \item \textbf{Wrapping:} Mylar on $\pm Y$ and $\pm Z$ faces;
            $R=\MylarR$, specular lobe $=\MylarSL$,
            $\sigma_\alpha=\MylarSA^\circ$.
      \item \textbf{Readout:} SUM4 leading-edge at 4\,PE;
            groups $\{0\text{--}3\}/\{4\text{--}7\}$ (left),
            $\{8\text{--}11\}/\{12\text{--}15\}$ (right).
      \item \textbf{Estimator:}
            $\sigma_\mathrm{single} = \sigma_{LR}/\sqrt{2}$.
      \item \textbf{Beam:} $\mu^-$, 1\,GeV, normal incidence.
    \end{itemize}
    \column{0.43\textwidth}
    \begin{block}{SiPM model}
      \textbf{\SiPMModel} (single-element, Broadcom).\\
      PDE \PDEpeak\ @ 420\,nm.\\
      \small\textit{Note: macro uses \texttt{\SiPMModelPDE} string
      only for the PDE curve; physical device is the single-element P014M.}
    \end{block}
    \begin{block}{Software stack}
      Geant4 11.4.0; SSLG4/OPSim.
    \end{block}
  \end{columns}
\end{frame}

% ─── FRAME 4: Physics inputs ──────────────────────────────────────────────────
\begin{frame}{Physics inputs}
  \begin{columns}[T]
    \column{0.50\textwidth}
    \begin{block}{EJ-204 (OPSC-101)}
      \small
      \begin{tabular}{lll}
        \toprule
        Property & Value & Unit \\
        \midrule
        Scint.\ yield & \ScintYield & \\
        Rise time $\tau_r$ & \TauRise & \\
        Decay time $\tau_d$ & \TauDecay & \\
        Bulk atten.\ $\lambda$ & \LamBulk & cm \\
        \bottomrule
      \end{tabular}
    \end{block}
    \begin{block}{Mylar wrapping ($\pm Y,\pm Z$)}
      \small $R=\MylarR$; specular lobe $=\MylarSL$;
      $\sigma_\alpha=\MylarSA^\circ$
    \end{block}
    \column{0.48\textwidth}
    \begin{block}{SiPM: \SiPMModel}
      \small
      \begin{tabular}{ll}
        \toprule
        Property & Value \\
        \midrule
        PDE (peak) & \PDEpeak\ @ 420\,nm \\
        Type & Single-element \\
        PDE curve model & \SiPMModelPDE \\
        \bottomrule
      \end{tabular}
    \end{block}
    \begin{block}{Scan parameters}
      \small \NPos\ positions; \NEv\ events each;
      $\mu^-$ 1\,GeV normal incidence.
    \end{block}
  \end{columns}
\end{frame}

% ─── FRAME 5: Photon count vs position ────────────────────────────────────────
\begin{frame}{Photon count at ends vs.\ beam position}
  \begin{columns}[T]
    \column{0.65\textwidth}
    \includegraphics[width=\textwidth]{figures/fig_npe_vs_x.png}
    \column{0.33\textwidth}
    \begin{block}{Key values}
      \small
      At $x=0$:\\
      $\langle N^\mathrm{L}_\mathrm{pe}\rangle = \NpeLeftZero$\,PE\\
      $\langle N^\mathrm{R}_\mathrm{pe}\rangle = \NpeRightZero$\,PE
      \medskip

      Near $x=-690$\,mm:\\
      Left $\approx \NpeLeftNear$\,PE\\
      Right $\approx \NpeRightNear$\,PE
    \end{block}
    \begin{alertblock}{\small Mylar role}
      \small Without Mylar the far-end count would collapse;
      Mylar recirculates photons and sustains far-end yield.
      Log scale shows exponential falloff.
    \end{alertblock}
  \end{columns}
\end{frame}

% ─── FRAME 6: Attenuation curve & λ_eff ──────────────────────────────────────
\begin{frame}{Attenuation curve \& effective attenuation length $\lambda_\mathrm{eff}$}
  \begin{columns}[T]
    \column{0.65\textwidth}
    \includegraphics[width=\textwidth]{figures/fig_attenuation.png}
    \column{0.33\textwidth}
    \begin{block}{Fit results}
      \small
      \begin{tabular}{ll}
        \toprule
        Side & $\lambda_\mathrm{eff}$ (cm) \\
        \midrule
        Left     & \LamEffL \\
        Right    & \LamEffR \\
        Combined & \LamEffC \\
        \bottomrule
      \end{tabular}
    \end{block}
    \begin{alertblock}{\small $\lambda_\mathrm{eff} \ll \lambda_\mathrm{bulk}$}
      \small $\lambda_\mathrm{eff}\approx\LamEffC$ vs.\
      bulk $\lambda=\LamBulk$.
      Mylar ($R=\MylarR$) recirculates light but each bounce
      has a $(1-R)$ loss. Many bounces $\Rightarrow$ short
      effective range.
    \end{alertblock}
  \end{columns}
\end{frame}

% ─── FRAME 7: Timing distributions ───────────────────────────────────────────
\begin{frame}{Arrival-time distributions: FPT at representative positions}
  \begin{columns}[T]
    \column{0.65\textwidth}
    \includegraphics[width=\textwidth]{figures/fig_timing_histograms.png}
    \column{0.33\textwidth}
    \begin{block}{Distributions shown}
      \small
      First Photon Time (FPT) at End-Left and End-Right per event
      (\textbf{not normalized}).
      Positions: $x\in\{-690,-400,0,+400,+690\}$\,mm.
      $N = \NEv$ events per position.
    \end{block}
    \begin{alertblock}{\small Key feature}
      \small Distribution shifts to larger $t$ and broadens as beam
      moves away from the respective end --- from longer propagation
      path and fewer photons (statistical minimum effect).
    \end{alertblock}
  \end{columns}
\end{frame}

% ─── FRAME 8: FPT vs t50 ─────────────────────────────────────────────────────
\begin{frame}{FPT vs.\ $t_{50}$: comparing timing estimators}
  \begin{columns}[T]
    \column{0.65\textwidth}
    \includegraphics[width=\textwidth]{figures/fig_fpt_vs_t50.png}
    \column{0.33\textwidth}
    \begin{block}{Estimators}
      \small
      \textbf{FPT} = first photon arrival time per event.\\
      $t_{50}$ = median photon arrival time per event.
    \end{block}
    \begin{block}{Observation}
      \small $t_{50}$ is shifted to later times relative to FPT,
      due to Mylar-reflected photons with longer paths.
      The shift grows with distance from the end
      (more reflected photons at large $d$).
    \end{block}
  \end{columns}
\end{frame}

% ─── FRAME 9: Mean arrival time & v_eff ──────────────────────────────────────
\begin{frame}{Mean arrival time \& effective velocity $v_\mathrm{eff}$}
  \begin{columns}[T]
    \column{0.65\textwidth}
    \includegraphics[width=\textwidth]{figures/fig_mean_arrival_time.png}
    \column{0.33\textwidth}
    \begin{block}{$v_\mathrm{eff}$ summary}
      \small
      \begin{tabular}{ll}
        \toprule
        Estimator & $v_\mathrm{eff}$ (mm/ns) \\
        \midrule
        FPT       & \VeffFPT \\
        $t_{50}$  & \VeffTfifty \\
        SUM4 $\Delta t$ & \VeffSUMfour \\
        \bottomrule
      \end{tabular}
    \end{block}
    \begin{alertblock}{\small $v_\mathrm{eff}$(FPT) anomaly}
      \small FPT bias: near the end ($N_\mathrm{pe}\sim 3000$)
      the statistical minimum is very small; far away
      ($N_\mathrm{pe}\sim 17$) it is larger.
      This $N$-dependent effect inflates the apparent slope,
      yielding $v_\mathrm{eff}(\mathrm{FPT}) > v_\mathrm{eff}(\Delta t)$.
      The SUM4 $\Delta t$ value (\VeffSUMfour) is more reliable.
    \end{alertblock}
  \end{columns}
\end{frame}

% ─── FRAME 10: σ_t(x) ────────────────────────────────────────────────────────
\begin{frame}{Intrinsic timing resolution $\sigma_t(x)$}
  \begin{columns}[T]
    \column{0.65\textwidth}
    \includegraphics[width=\textwidth]{figures/fig_sigma_t.png}
    \column{0.33\textwidth}
    \begin{block}{$\sigma_t$ at key positions}
      \small
      \begin{tabular}{rl}
        \toprule
        $x$ (mm) & $\sigma_\mathrm{single}$ (ps) \\
        \midrule
        $-690$ & \SigmaAtMsixnine \\
        $-400$ & \SigmaAtMfour    \\
        $0$    & \SigmaAtZero     \\
        $+400$ & \SigmaAtFour     \\
        $+690$ & \SigmaAtSixnine  \\
        \bottomrule
      \end{tabular}
    \end{block}
    \begin{alertblock}{\small Intrinsic only}
      \small $\sigma_t$ shown = optical + SUM4 estimator.
      SiPM SPTR and electronics jitter \textbf{not} included.
      $\sigma_t^\mathrm{total} = \sqrt{\sigma_t^2
      + \sigma_\mathrm{SPTR}^2 + \sigma_\mathrm{elec}^2}$.
    \end{alertblock}
  \end{columns}
\end{frame}

% ─── FRAME 11: Photon budget ──────────────────────────────────────────────────
\begin{frame}{Photon budget: fate of scintillation photons}
  \begin{columns}[T]
    \column{0.65\textwidth}
    \includegraphics[width=\textwidth]{figures/fig_photon_budget.png}
    \column{0.33\textwidth}
    \begin{block}{Budget (all 31 positions)}
      \small
      \begin{tabular}{ll}
        \toprule
        Fate & Fraction \\
        \midrule
        Escaped       & \BudgetEsc \\
        Enter SiPM    & \BudgetSiPM \\
        Detected      & \BudgetDet \\
        \bottomrule
      \end{tabular}
    \end{block}
    \begin{alertblock}{\small Mylar absorption dominates}
      \small $\approx98\%$ absorbed by Mylar over many bounces
      ($R=\MylarR$ per reflection).
      ``Reflector panel crossings'' counts every bounce,
      not unique photons.
    \end{alertblock}
  \end{columns}
\end{frame}

% ─── FRAME 12: Summary & next steps ──────────────────────────────────────────
\begin{frame}{Summary \& next steps}
  \begin{columns}[T]
    \column{0.55\textwidth}
    \begin{block}{Summary}
      \begin{itemize}
        \item End-only + Mylar = clean simulation limit.
        \item $\lambda_\mathrm{eff} \approx \LamEffC$
              $\ll \lambda_\mathrm{bulk} = \LamBulk$.
              Mylar absorption ($1-R$ per bounce) explains reduction.
        \item $\sigma_t^\mathrm{intr.}$: \SigmaAtZero\ (center)
              to \SigmaAtMsixnine\ (ends).
        \item $v_\mathrm{eff}(\Delta t) \approx \VeffSUMfour$.
        \item Photon yield: Mylar absorbs $\gtrsim 98\%$.
      \end{itemize}
    \end{block}
    \column{0.43\textwidth}
    \begin{block}{Next steps}
      \begin{enumerate}
        \item Fold in electronic jitter:
              $\sigma_\mathrm{SPTR}=106$\,ps
              $\oplus\,\sigma_\mathrm{FastIC}=10$\,ps.
        \item Leading-edge time-walk correction.
        \item Compare with April-2026 TB data (different config:
              top readout + Tyvek).
        \item Scan Mylar $R\in[0.85,0.95]$ for systematic study.
      \end{enumerate}
    \end{block}
  \end{columns}
\end{frame}

% ════════════════════════════════════════════════════════════════════════════════
% APPENDIX
% ════════════════════════════════════════════════════════════════════════════════
\appendix

\begin{frame}{Appendix: simulation parameters}
  \small
  \begin{tabular}{llll}\toprule
Parameter & Value & Unit & Source \\
\midrule
Bar material & EJ-204 (OPSC-101) &  & SSLG4 \\
Bar dimensions & 1400\,$\times$\,60\,$\times$\,10 & mm$^3$ & GDML \\
Readout & End-only ($\pm X$ faces) &  & geometry \\
Wrapping ($\pm Y,\pm Z$) & Mylar &  & run\_metadata \\
Mylar reflectivity & 0.90 &  & run\_metadata \\
Mylar specular lobe & 1.0 &  & run\_metadata \\
Mylar $\sigma_\alpha$ & 0.1 & deg & run\_metadata \\
Scint.~yield & 10000 & $\gamma$/MeV & SSLG4 (EJ-204) \\
Rise time $\tau_r$ & 0.7 & ns & SSLG4 (EJ-204) \\
Decay time $\tau_d$ & 1.8 & ns & SSLG4 (EJ-204) \\
Bulk att.~$\lambda$ & 160 & cm & EJ-204 datasheet \\
SiPM model & AFBR-S4N66P014M &  & hardware \\
SiPM PDE (peak) & 63 & \% @ 420\,nm & datasheet \\
Beam & $\mu^-$, 1\,GeV &  & macro \\
Positions & 31 &  & run\_metadata \\
Events/position & 2000 &  & run\_metadata \\
Geant4 version & 11.4.0 &  & log \\
\bottomrule\end{tabular}

\end{frame}

\begin{frame}{Appendix: $\sigma_t$ per-position summary}
  \centering\small
  \begin{tabular}{rrrrl}\toprule
$x$ (mm) & $N_\mathrm{trig}$ & Trig.~eff. & $\sigma_\mathrm{single}$ (ps) & Err. (ps) \\
\midrule
$-690$ & 1911 & 0.956 & 406.2 & 26.6 \\
$-400$ & 2000 & 1.000 & 210.6 & 11.2 \\
$+0$ & 2000 & 1.000 & 141.5 & 4.3 \\
$+400$ & 2000 & 1.000 & 209.3 & 8.9 \\
$+690$ & 1911 & 0.956 & 372.1 & 19.1 \\
\bottomrule\end{tabular}

  \par\medskip
  \small All values in ps.
  $\sigma_\mathrm{single}=\sigma_{LR}/\sqrt{2}$ from SUM4 leading-edge
  core Gaussian fit.
\end{frame}

\begin{frame}{Appendix: attenuation fit details}
  \centering\small
  \begin{tabular}{lrrrrl}\toprule
Side & $N_0$ (PE) & $\sigma(N_0)$ & $\lambda_\mathrm{eff}$ (cm) & $\sigma(\lambda)$ & $N$ \\
\midrule
Left & 811.5 & $\pm$2.2 & 32.16 & $\pm$0.04 & 31 \\
Right & 853.1 & $\pm$2.2 & 31.41 & $\pm$0.04 & 31 \\
Combined & 832.5 & $\pm$1.6 & 31.78 & $\pm$0.03 & 62 \\
\bottomrule\end{tabular}

  \par\medskip
  \small Fit: $N_\mathrm{pe}(d)=N_0\exp(-d/\lambda_\mathrm{eff})$.
  31 points per side; 62 combined.
  Errors are $1\sigma$ from covariance matrix.
\end{frame}

\begin{frame}{Appendix: method \& caveats}
  \small
  \begin{itemize}
    \item \textbf{Geometry:} Bar $1400\times60\times10$\,mm$^3$; long axis $X$.
          End SiPMs on $\pm X$; Mylar on $\pm Y,\pm Z$.
    \item \textbf{SUM4 trigger:} leading-edge at 4\,PE.
          Groups $\{0\text{--}3\}/\{4\text{--}7\}$ (left);
          $\{8\text{--}11\}/\{12\text{--}15\}$ (right).
    \item \textbf{$\sigma_t$ is intrinsic only:} optical + estimator;
          SPTR ($\approx106$\,ps) and FastIC ($\approx10$\,ps)
          not folded in.
    \item \textbf{SiPM naming:} macro string \texttt{\SiPMModelPDE}
          is used \emph{only} for the PDE curve;
          physical device is the single-element \textbf{\SiPMModel}.
    \item \textbf{No TB counterpart:} April-2026 TB used top readout
          + Tyvek wrapping $\neq$ end-only + Mylar.
    \item \textbf{Longitudinal axis $= X$} (not $Z$ or $Y$).
  \end{itemize}
\end{frame}

\end{document}
